\documentclass[10pt,a4paper]{amsart}
\usepackage{amsmath,amssymb,amsthm,graphicx}
% For dealing with references we use the comment environment
\usepackage{verbatim}
\newenvironment{reference}{\comment}{\endcomment}
%\newenvironment{reference}{}{}
\newenvironment{slogan}{\comment}{\endcomment}
\newenvironment{history}{\comment}{\endcomment}
% For commutative diagrams we use Xy-pic
\usepackage[all]{xy}

% We use 2cell for 2-commutative diagrams.
\xyoption{2cell}
\UseAllTwocells
\theoremstyle{plain}
\newtheorem{theorem}[subsection]{Theorem}
\newtheorem{proposition}[subsection]{Proposition}
\newtheorem{lemma}[subsection]{Lemma}

\theoremstyle{definition}
\newtheorem{definition}[subsection]{Definition}
\newtheorem{example}[subsection]{Example}
\newtheorem{exercise}[subsection]{Exercise}
\newtheorem{situation}[subsection]{Situation}

\theoremstyle{remark}
\newtheorem{remark}[subsection]{Remark}
\newtheorem{remarks}[subsection]{Remarks}

\numberwithin{equation}{subsection}
% Macros
%
\def\lim{\mathop{\mathrm{lim}}\nolimits}
\def\colim{\mathop{\mathrm{colim}}\nolimits}
\def\Spec{\mathop{\mathrm{Spec}}}
\def\Hom{\mathop{\mathrm{Hom}}\nolimits}
\def\Ext{\mathop{\mathrm{Ext}}\nolimits}
\def\SheafHom{\mathop{\mathcal{H}\!\mathit{om}}\nolimits}
\def\SheafExt{\mathop{\mathcal{E}\!\mathit{xt}}\nolimits}
\def\Sch{\mathit{Sch}}
\def\Mor{\mathop{\mathrm{Mor}}\nolimits}
\def\Ob{\mathop{\mathrm{Ob}}\nolimits}
\def\Sh{\mathop{\mathit{Sh}}\nolimits}
\def\NL{\mathop{N\!L}\nolimits}
\def\CH{\mathop{\mathrm{CH}}\nolimits}
\def\proetale{{pro\text{-}\acute{e}tale}}
\def\etale{{\acute{e}tale}}
\def\QCoh{\mathit{QCoh}}
\def\Ker{\mathop{\mathrm{Ker}}}
\def\Im{\mathop{\mathrm{Im}}}
\def\Coker{\mathop{\mathrm{Coker}}}
\def\Coim{\mathop{\mathrm{Coim}}}

% Boxtimes
%
\DeclareMathSymbol{\boxtimes}{\mathbin}{AMSa}{"02}

%
% Macros for moduli stacks/spaces
%
\def\QCohstack{\mathcal{QC}\!\mathit{oh}}
\def\Cohstack{\mathcal{C}\!\mathit{oh}}
\def\Spacesstack{\mathcal{S}\!\mathit{paces}}
\def\Quotfunctor{\mathrm{Quot}}
\def\Hilbfunctor{\mathrm{Hilb}}
\def\Curvesstack{\mathcal{C}\!\mathit{urves}}
\def\Polarizedstack{\mathcal{P}\!\mathit{olarized}}
\def\Complexesstack{\mathcal{C}\!\mathit{omplexes}}
% \Pic is the operator that assigns to X its picard group, usage \Pic(X)
% \Picardstack_{X/B} denotes the Picard stack of X over B
% \Picardfunctor_{X/B} denotes the Picard functor of X over B
\def\Pic{\mathop{\mathrm{Pic}}\nolimits}
\def\Picardstack{\mathcal{P}\!\mathit{ic}}
\def\Picardfunctor{\mathrm{Pic}}
\def\Deformationcategory{\mathcal{D}\!\mathit{ef}}
\usepackage{cleveref}
\crefname{originaldefinition}{original definition}{original definitions}
\Crefname{originaldefinition}{Original definition}{Original definitions}
\makeatletter
\newlabel{algebra-definition-analytically-unramified}{{032X}{source}{}{sourcepointer-1}{}}
\newlabel{algebra-definition-analytically-unramified@cref}{{[originaldefinition][0][]032X}{[1][0][]source}{}{}{}}
\makeatother
\crefname{originallemma}{original lemma}{original lemmas}
\Crefname{originallemma}{Original lemma}{Original lemmas}
\makeatletter
\newlabel{algebra-lemma-analytically-unramified-easy}{{032Y}{source}{}{sourcepointer-2}{}}
\newlabel{algebra-lemma-analytically-unramified-easy@cref}{{[originallemma][0][]032Y}{[1][0][]source}{}{}{}}
\makeatother
\crefname{originallemma}{original lemma}{original lemmas}
\Crefname{originallemma}{Original lemma}{Original lemmas}
\makeatletter
\newlabel{descent-lemma-reduced-local-smooth}{{034E}{source}{}{sourcepointer-3}{}}
\newlabel{descent-lemma-reduced-local-smooth@cref}{{[originallemma][0][]034E}{[1][0][]source}{}{}{}}
\makeatother
\crefname{originallemma}{original lemma}{original lemmas}
\Crefname{originallemma}{Original lemma}{Original lemmas}
\makeatletter
\newlabel{algebra-lemma-completion-faithfully-flat}{{00MC}{source}{}{sourcepointer-4}{}}
\newlabel{algebra-lemma-completion-faithfully-flat@cref}{{[originallemma][0][]00MC}{[1][0][]source}{}{}{}}
\makeatother
\crefname{originalsection}{original section}{original sections}
\Crefname{originalsection}{Original section}{Original sections}
\makeatletter
\newlabel{algebra-section-going-up}{{00HU}{source}{}{sourcepointer-5}{}}
\newlabel{algebra-section-going-up@cref}{{[originalsection][0][]00HU}{[1][0][]source}{}{}{}}
\makeatother
\crefname{originallemma}{original lemma}{original lemmas}
\Crefname{originallemma}{Original lemma}{Original lemmas}
\makeatletter
\newlabel{algebra-lemma-generically-finite}{{02ML}{source}{}{sourcepointer-6}{}}
\newlabel{algebra-lemma-generically-finite@cref}{{[originallemma][0][]02ML}{[1][0][]source}{}{}{}}
\makeatother
\crefname{originalsection}{original section}{original sections}
\Crefname{originalsection}{Original section}{Original sections}
\makeatletter
\newlabel{morphisms-section-generically-finite}{{02NV}{source}{}{sourcepointer-7}{}}
\newlabel{morphisms-section-generically-finite@cref}{{[originalsection][0][]02NV}{[1][0][]source}{}{}{}}
\makeatother
\crefname{originallemma}{original lemma}{original lemmas}
\Crefname{originallemma}{Original lemma}{Original lemmas}
\makeatletter
\newlabel{algebra-lemma-Noetherian-complete-local-Nagata}{{032W}{source}{}{sourcepointer-8}{}}
\newlabel{algebra-lemma-Noetherian-complete-local-Nagata@cref}{{[originallemma][0][]032W}{[1][0][]source}{}{}{}}
\makeatother
\crefname{originallemma}{original lemma}{original lemmas}
\Crefname{originallemma}{Original lemma}{Original lemmas}
\makeatletter
\newlabel{algebra-lemma-nagata-in-reduced-finite-type-finite-integral-closure}{{03GH}{source}{}{sourcepointer-9}{}}
\newlabel{algebra-lemma-nagata-in-reduced-finite-type-finite-integral-closure@cref}{{[originallemma][0][]03GH}{[1][0][]source}{}{}{}}
\makeatother
\crefname{originallemma}{original lemma}{original lemmas}
\Crefname{originallemma}{Original lemma}{Original lemmas}
\makeatletter
\newlabel{algebra-lemma-descend-properties-modules}{{03C4}{source}{}{sourcepointer-10}{}}
\newlabel{algebra-lemma-descend-properties-modules@cref}{{[originallemma][0][]03C4}{[1][0][]source}{}{}{}}
\makeatother
\title{Lemma formally unramified in resolve.tex}
\author{The Stacks Project Authors}
\date{stacks-project:a04446e57ec1fbc252a871afcec7752fb2807b14; GNU Free Documentation License 1.2 (as published in the Stacks Project COPYING)}
\begin{document}
\maketitle
\noindent\textit{Editorial note.} All mathematical statement and proof text is reproduced verbatim from the author's chapter resolve.tex of the Stacks Project, fixed at commit a04446e57ec1fbc252a871afcec7752fb2807b14. The pinned preamble is supplied as its own file; the wrapper class, the xr-hyper wiring and the \textbackslash{}externaldocument list are omitted because the assembler supplies the class.\par\medskip
\noindent\textit{Editorial note.} Numbering inside this standalone excerpt is editorial, not the full Stacks book print numbering.\par\medskip
\noindent\textit{Editorial note.} The candidate is from the corrected band-I follow-on screen (line C) of the extractor's stacks-candidates line: its author label is absent from the live master in every form searched and it belongs to none of the earlier judged pools. The H2 judgement and its reason were made in this module by reading the author's own statement and complete proof in the fixed checkout. Original candidate key: stacks:resolve.tex:lemma-formally-unramified.\par\medskip
\noindent\textit{Editorial note.} Reference labels in the proof resolve to navigation pointers to the original author source; these pointers do not reproduce proofs of the prerequisites.\par\medskip
\noindent\textit{Allowed standard prerequisites (editorial).}\begin{itemize}
\item Prerequisites named by the author inside the counted span. The counted span is reproduced verbatim, including the author's own citations. Its prerequisites are the earlier results of the same fixed Stacks Project version that the author cites there and in the chapter's lead-in.
\end{itemize}\medskip
\ifdefined\hypertarget\hypertarget{sourcepointer-1}{}\fi
\noindent\textit{Source reference.} Original definition 032X: sources/algebra.tex, source lines 45781--45781. Navigation pointer used in the counted span; the span contains the pointer only and the cited result is not reproduced inside it. Registered as a non-load-bearing navigation stub in the established form for Stacks material. No bibliography citation is involved.\par\medskip
\ifdefined\hypertarget\hypertarget{sourcepointer-2}{}\fi
\noindent\textit{Source reference.} Original lemma 032Y: sources/algebra.tex, source lines 45808--45808. Navigation pointer used in the counted span; the span contains the pointer only and the cited result is not reproduced inside it. Registered as a non-load-bearing navigation stub in the established form for Stacks material. No bibliography citation is involved.\par\medskip
\ifdefined\hypertarget\hypertarget{sourcepointer-3}{}\fi
\noindent\textit{Source reference.} Original lemma 034E: sources/descent.tex, source lines 4952--4952. Navigation pointer used in the counted span; the span contains the pointer only and the cited result is not reproduced inside it. Registered as a non-load-bearing navigation stub in the established form for Stacks material. No bibliography citation is involved.\par\medskip
\ifdefined\hypertarget\hypertarget{sourcepointer-4}{}\fi
\noindent\textit{Source reference.} Original lemma 00MC: sources/algebra.tex, source lines 22977--22977. Navigation pointer used in the counted span; the span contains the pointer only and the cited result is not reproduced inside it. Registered as a non-load-bearing navigation stub in the established form for Stacks material. No bibliography citation is involved.\par\medskip
\ifdefined\hypertarget\hypertarget{sourcepointer-5}{}\fi
\noindent\textit{Source reference.} Original section 00HU: sources/algebra.tex, source lines 9515--9515. Navigation pointer used in the counted span; the span contains the pointer only and the cited result is not reproduced inside it. Registered as a non-load-bearing navigation stub in the established form for Stacks material. No bibliography citation is involved.\par\medskip
\ifdefined\hypertarget\hypertarget{sourcepointer-6}{}\fi
\noindent\textit{Source reference.} Original lemma 02ML: sources/algebra.tex, source lines 30418--30418. Navigation pointer used in the counted span; the span contains the pointer only and the cited result is not reproduced inside it. Registered as a non-load-bearing navigation stub in the established form for Stacks material. No bibliography citation is involved.\par\medskip
\ifdefined\hypertarget\hypertarget{sourcepointer-7}{}\fi
\noindent\textit{Source reference.} Original section 02NV: sources/morphisms.tex, source lines 13184--13184. Navigation pointer used in the counted span; the span contains the pointer only and the cited result is not reproduced inside it. Registered as a non-load-bearing navigation stub in the established form for Stacks material. No bibliography citation is involved.\par\medskip
\ifdefined\hypertarget\hypertarget{sourcepointer-8}{}\fi
\noindent\textit{Source reference.} Original lemma 032W: sources/algebra.tex, source lines 45755--45755. Navigation pointer used in the counted span; the span contains the pointer only and the cited result is not reproduced inside it. Registered as a non-load-bearing navigation stub in the established form for Stacks material. No bibliography citation is involved.\par\medskip
\ifdefined\hypertarget\hypertarget{sourcepointer-9}{}\fi
\noindent\textit{Source reference.} Original lemma 03GH: sources/algebra.tex, source lines 45640--45640. Navigation pointer used in the counted span; the span contains the pointer only and the cited result is not reproduced inside it. Registered as a non-load-bearing navigation stub in the established form for Stacks material. No bibliography citation is involved.\par\medskip
\ifdefined\hypertarget\hypertarget{sourcepointer-10}{}\fi
\noindent\textit{Source reference.} Original lemma 03C4: sources/algebra.tex, source lines 20307--20307. Navigation pointer used in the counted span; the span contains the pointer only and the cited result is not reproduced inside it. Registered as a non-load-bearing navigation stub in the established form for Stacks material. No bibliography citation is involved.\par\medskip
\begin{lemma}
\label{lemma-formally-unramified}
Let $(A, \mathfrak m)$ be a local Noetherian ring. Let
$X$ be a scheme over $A$. Assume
\begin{enumerate}
\item $A$ is analytically unramified
(Algebra, Definition \ref{algebra-definition-analytically-unramified}),
\item $X$ is locally of finite type over $A$, and
\item $X \to \Spec(A)$ is \'etale at the generic points of irreducible
components of $X$.
\end{enumerate}
Then the normalization of $X$ is finite over $X$.
\end{lemma}

\begin{proof}
Since $A$ is analytically unramified it is reduced
by Algebra, Lemma \ref{algebra-lemma-analytically-unramified-easy}.
Since the normalization of $X$ depends only on the reduction
of $X$, we may replace $X$ by its reduction $X_{red}$; note
that $X_{red} \to X$ is an isomorphism over the open $U$ where
$X \to \Spec(A)$ is \'etale because $U$ is reduced
(Descent, Lemma \ref{descent-lemma-reduced-local-smooth})
hence condition (3) remains true after this replacement.
In addition we may and do assume that $X = \Spec(B)$ is affine.

\medskip\noindent
The map
$$
K = \prod\nolimits_{\mathfrak p \subset A\text{ minimal}} \kappa(\mathfrak p)
\longrightarrow
K^\wedge = \prod\nolimits_{\mathfrak p^\wedge \subset A^\wedge\text{ minimal}}
\kappa(\mathfrak p^\wedge)
$$
is injective because $A \to A^\wedge$ is faithfully flat
(Algebra, Lemma \ref{algebra-lemma-completion-faithfully-flat})
hence induces a surjective map between sets of minimal primes
(by going down for flat ring maps, see
Algebra, Section \ref{algebra-section-going-up}).
Both sides are finite products of fields as our rings are Noetherian.
Let $L = \prod_{\mathfrak q \subset B\text{ minimal}} \kappa(\mathfrak q)$.
Our assumption (3) implies that $L = B \otimes_A K$ and that
$K \to L$ is a finite \'etale ring map (this is true
because $A \to B$ is generically finite, for example use
Algebra, Lemma \ref{algebra-lemma-generically-finite}
or the more detailed results in Morphisms, Section
\ref{morphisms-section-generically-finite}).
Since $B$ is reduced we see that $B \subset L$.
This implies that
$$
C = B \otimes_A A^\wedge \subset
L \otimes_A A^\wedge = L \otimes_K K^\wedge = M
$$
Then $M$ is the total ring of fractions of $C$ and
is a finite product of fields as a finite separable
algebra over $K^\wedge$. It follows that $C$ is reduced
and that its normalization $C'$ is the integral closure of
$C$ in $M$. The normalization $B'$ of $B$ is the integral
closure of $B$ in $L$. By flatness of $A \to A^\wedge$
we obtain an injective map $B' \otimes_A A^\wedge \to M$ whose
image is contained in $C'$. Picture
$$
B' \otimes_A A^\wedge \longrightarrow C'
$$
As $A^\wedge$ is Nagata (by
Algebra, Lemma \ref{algebra-lemma-Noetherian-complete-local-Nagata}),
we see that $C'$ is finite over
$C = B \otimes_A A^\wedge$ (see
Algebra, Lemmas
\ref{algebra-lemma-Noetherian-complete-local-Nagata} and
\ref{algebra-lemma-nagata-in-reduced-finite-type-finite-integral-closure}).
As $C$ is Noetherian, we conclude that
$B' \otimes_A A^\wedge$ is finite over $C = B \otimes_A A^\wedge$.
Therefore by faithfully flat descent
(Algebra, Lemma \ref{algebra-lemma-descend-properties-modules})
we see that $B'$ is finite over $B$ which is what we had to show.
\end{proof}

\end{document}
